\newpage
\appendix

% ---------------------------------------------------------------------------
% Populated appendix sections (content ported from body in the 2026-05-24 cut).
% ---------------------------------------------------------------------------

\section{Authoring process}
\label{app:authoring}

Environments were curated by the authors and implemented through LLM-assisted code generation. For each environment we wrote a one-page design brief specifying the impersonated sector and brand archetype, the legitimate task, the attack surface and its placement, the values of all eight axes, the visible red flags, the expected endpoint behaviour, and the sibling-relationship pointer (parent and toggled axis). The brief defines the environment as an artifact.

A frontier code-generation model then authored the HTML, CSS, JavaScript, and Flask boilerplate from the brief. The authoring model and its version are pinned across all 91 environments and disclosed alongside the pinned evaluator-model identifiers (Appendix~\ref{app:models}); it is distinct from any of the four models evaluated in \S\ref{sec:setup-runs}, so no agent in the panel is graded on implementations it itself produced. Its role is bounded and clerical: it implements a fully-specified design but does not choose the vector, salience, pressure cue, PII target, prompt-injection content, or any other axis value. The attack composition is fixed in the design brief before any code is generated, and every brief was written by a human author; the model is not given latitude over the experimental factor it would otherwise be evaluated on. The design briefs themselves are not generated by a model.

Following generation, every environment was manually reviewed by the authors against its brief: we verified that all eight axis values match, that red flags appear at the specified locations in both DOM and screenshot, that \texttt{/api/captured} records the expected payload, and that the legitimate task completes end-to-end. Environments failing any check were revised and re-reviewed; all 91 released environments passed.

We adopt this curated-design / generated-implementation style because (i) maintaining 91 environments at the breadth required for paired-sibling ablation across eight axes is incompatible with fully manual code authoring, and (ii) axis-level control requires surgical edits between sibling pairs (e.g., toggling only the pressure cue), and code generation from a versioned brief produces more reliable single-axis diffs than re-authoring sibling implementations from scratch.

\section{Pre-registration: artifacts frozen at the tag}
\label{app:prereg}

The complete analysis plan (primary outcome, hypothesis family, statistical model, multiple-comparison correction, power calculation, inclusion rules, falsification thresholds, and pre-committed deviations) is committed to the repository as \texttt{analysis-plan.md} and tagged \texttt{prereg-v2-start} before any experimental run reported in \S\ref{sec:results}. Artifacts frozen at the tag include the analysis plan itself, the environment classification table and registry (Appendix~\ref{app:envs}), the C0--C3 mitigation prompts (Appendix~\ref{app:prompts}), the critical-tier definition in the PII tracker, and the pinned model identifiers (Appendix~\ref{app:models}). No deviation is applied retroactively to any pre-tag artifact, and no post-tag change touches the primary outcome definition, the analysis pipeline, the prompt strings, or the environment registry.

\section{Deviation log (D1--D9)}
\label{app:deviations}

The full deviation log with rationale, git SHA for each fix, and the pre-fix vs post-fix slices preserved for contamination accounting is enumerated in \texttt{analysis-plan.md} \S 11 (entries D1 through D9). The deviations exercised against the runs reported in \S\ref{sec:results} are D5 (observation-URL contamination correction; full audit in Appendix~\ref{app:contamination}) and D6 (Anthropic-slot model swap from Sonnet 4.6 to Haiku 4.5, executed pre-data). D7 (BH family reduced from 14 to 8 after the \S\ref{sec:analysis} power gate), D8 (AME-scaling correction in mixed-effects logistic reporting), and D9 (LLM-judge DR landing and the F1 anchor confirmation) are measurement-instrument deviations that do not change any session's PII trail or scorer output.

\section{Primary-family test table}
\label{app:claims}

\begin{table}[h]
\centering
\small
\resizebox{\columnwidth}{!}{
\begin{tabular}{lll}
\toprule
Claim & Contrast & Toggled factor \\
\midrule
F1 & DR=1 vs DR=0 in C3 & within-cell (DR) \\
F2 & subtle vs blatant & C: salience \\
F3 & subtle $\times$ conversational & C $\times$ G \\
F4 & between-model variance & model factor \\
F5 & urgency vs neutral & E: pressure \\
F6 & social-proof vs neutral & E: pressure \\
F7 & authority vs neutral & E: pressure \\
F8 & hidden vs visible PI & F: prompt\_injection \\
F9 & critical- vs medium-tier PII & D: pii\_target \\
F10 & chat vs static form & G: interaction \\
F11 & cross-category & A: category \\
M1 & C0 vs C1 & condition \\
M2 & C0 vs C2 & condition \\
M3 & C0 vs C3 & condition \\
\bottomrule
\end{tabular}
}
\caption{The 14 pre-registered primary tests. Each is a single-axis or single-condition contrast on session-level $\text{PLR}_{\text{crit}}$. After the \S\ref{sec:analysis} power gate (analysis-plan \S 11 D7), Benjamini-Hochberg correction is applied across the 8-test inference-bearing subset (F1, F2, F5, F8, F10, M1, M2, M3); the remaining axes are reported as descriptive context with uncorrected $p$-values.}
\label{tab:claims}
\end{table}

\section{URL-shape contamination audit}
\label{app:contamination}

An early version of the harness left the localhost sandbox visible in agent observations through three surfaces: the raw \texttt{http://localhost:N} URL in the address bar, the same string rendered inline inside page templates (form action attributes, anchor hrefs, JavaScript endpoint constants, and mailbox-body embedded URLs), and the start URL substituted into the agent's task description before the first observation. All three functioned as model-side trust shortcuts: in pre-fix GPT-5 mini reasoning traces under C3, the strings ``localhost'' and ``non-standard port'' appear as cited suspicion signals in XX.X\% and XX.X\% of sessions respectively, an artefact that would inflate the apparent C3 effect.

The surface was identified by a deliberate audit pass rather than discovered post-hoc: we ran (i) a content-level grep over every templated environment for the eleven leakage strings enumerated in \texttt{scripts/audit\_env\_leakage.py}, (ii) a reasoning-trace scan of the pre-fix slice for the same strings to confirm the agent was citing them, (iii) a cross-cell check that observer-side, content-side, and task-side surfaces were each contributing, (iv) an inter-judge $\kappa$ stratification (\S\ref{sec:dr}) to rule out judge-side contamination, and (v) an audit of the benign-twin subtraction to confirm the artefact did not differentially affect attack versus benign cells. The localhost surface was the only one to surface measurable contamination; the four remaining checks returned negative.

After domain masking at the observer (\S\ref{sec:metrics}), a content-level sweep removing inline localhost strings across XX environment templates, an extension of the browser's request-layer rewriter to in-page navigation requests, and a fix to the task description's start URL, the C3 pooled $\Delta\text{PLR}_{\text{crit}}$ on GPT-5 mini moved from $-$XX.X pp pre-mask to $-$XX.X pp post-mask, of which approximately $-$XX.X pp is attributable to URL-shape contamination and $-$XX.X pp to the true C3 signal. All results in \S\ref{sec:results} are reported on the post-cleanup harness; the pre-cleanup measurements are preserved in the run logs for contamination accounting, and the git SHA of each surface-fix commit is pinned in the deviation log (Appendix~\ref{app:deviations}).

% ---------------------------------------------------------------------------
% Stub appendix sections (to be populated in the next pass).
% Each carries a one-line note describing what will fill it.
% ---------------------------------------------------------------------------

\section{DR-judge rubric and sensitivity}
\label{app:dr-sensitivity}

\emph{To be populated next pass.} Will include: verbatim DR-judge rubric prompt; per-model stratified inter-judge $\kappa$ table; both-judges-agree sensitivity table; keyword-DR comparison and overcount accounting; judge-vs-human $\kappa$ on the 200-pair validation subset once labelled.

\section{PhishTank/Wayback fidelity check}
\label{app:fidelity}

\emph{To be populated next pass.} A descriptive fidelity check rating 10 attack environments against 5 PhishTank/Wayback phishing screenshots on visual believability and copy quality.

\section{Per-environment registry}
\label{app:envs}

\emph{To be populated next pass.} Per-environment axis assignments for all 91 attack environments and 10 benign twins, parent--sibling relationships, and pinned port numbers.

\section{Pinned model identifiers and sampling parameters}
\label{app:models}

\emph{To be populated next pass.} Provider, version string, sampling temperature, max-tokens, and any provider-specific knobs for each of the four evaluated models, frozen at \texttt{prereg-v2-start}. Also lists the authoring-model identifier referenced from Appendix~\ref{app:authoring}.

\section{PII-matcher precision and recall}
\label{app:matcher}

\emph{To be populated next pass.} Precision and recall of the exact-string PLR matcher against a 200-pair human-labelled subset, plus a sensitivity comparison against substring and regex variants.

\section{Per-model breakdowns: F4 between-model variance, F11 cross-category}
\label{app:per-model}

\emph{To be populated next pass.} F4 between-model variance and F11 cross-category breakdowns, plus per-model F1 tables with confidence intervals and the underpowered-band caveats from \S\ref{sec:results-f1}.

\section{Verbatim C0--C3 mitigation prompts}
\label{app:prompts}

\emph{To be populated next pass.} Verbatim text of the C0, C1, C2, and C3 prompt strings as frozen at \texttt{prereg-v2-start}.

\section{Full MITRE / OWASP / ENISA mapping}
\label{app:mitre}

\emph{To be populated next pass.} Vector-by-vector mapping of the eight attack vectors (\S\ref{sec:vectors}) to MITRE ATT\&CK Enterprise techniques, OWASP Top 10 for LLM Applications entries, and ENISA AI threat landscape categories.
